\chapter*{Introducción}
\addcontentsline{toc}{chapter}{Introducción}

\noindent
El presente informe documenta el desarrollo del Proyecto Integrador de la cátedra \textit{Teoría de la Información y la Comunicación}, cuyo objetivo consiste en diseñar e implementar un sistema de escritorio capaz de \textbf{compactar} y \textbf{proteger} archivos de texto mediante la aplicación de los algoritmos de Huffman y Hamming, respectivamente. El sistema permite al usuario comprimir un archivo utilizando el método de Huffman, protegerlo contra errores de transmisión mediante el código de Hamming (con módulos de 8, 1024 y 16384~bits), simular la introducción de errores aleatorios y, finalmente, recuperar el archivo original realizando los procesos inversos de desprotección y descompactación.

El desarrollo se estructuró de manera incremental, resolviendo primero cada laboratorio de forma aislada (Laboratorio~1: Hamming; Laboratorio~2: Huffman) para luego integrarlos en un único proyecto funcional (Laboratorio~3), acompañado de una interfaz gráfica moderna.

A lo largo de este informe se abordan los siguientes ejes: en primer lugar, se presentan las definiciones teóricas de los códigos de Hamming y Huffman; luego se detallan las decisiones tecnológicas y algorítmicas tomadas durante la programación; a continuación se ofrece un manual de usuario con las instrucciones de ejecución del programa; se describen los principales problemas encontrados y sus soluciones; y, finalmente, se exponen los resultados estadísticos obtenidos tras someter al sistema a distintas corridas con archivos de diferentes tamaños y contenidos.